\section{Conclusion}

Paper I establishes recursive research as a governed epistemic-mutation problem rather than a synonym for repeated model calls. ORION separates knowledge, relevance/search-universe, and governed method state; diagnoses typed residuals so that responsibility constrains authority to rewrite a formulation coordinate; and binds each accepted change to the closure that actually depends on that coordinate. Strong donor mechanisms for audit, diagnosis, repair, rollback, certificate enforcement, and mutation control are retained inside one protected scientific-escalation contract: lower-level alternatives must be excluded, a high-level change must produce a same-task counterfactual benefit, protected siblings must remain valid, and dependency impact must be explicit before a K/W/M mutation is admitted.

The historical negative result is part of that scientific story, but not as failed performance to be hidden or rebranded. The old 48-case broad-superiority hypothesis remains unsupported. Together with the frozen local falsifier, it identifies the unsafe rule the successor replaces: \emph{failure itself is not authority for high-level reformulation}. The publication-authorizing successor was therefore frozen prospectively around a stricter necessity condition rather than tuned by relabeling the old outcome. This creates a falsification-to-repair result: an over-broad escalation rule is rejected, its missing authority condition is isolated, and the repaired rule is then tested on fresh protected evidence.

That corrected contract produces a large and replicated result. Across the prospectively frozen v2.2.4 primary run and disjoint replication, ORION attains \textbf{1.000 protected hidden-shift success} while the strongest frozen parent reaches 0.494 and 0.483, respectively. The paired advantages exceed \textbf{0.50 in both runs} and clear the registered 0.10 practical margin, while ORION simultaneously maintains \textbf{zero unnecessary high-level reframes, zero protected-sibling regressions, and zero forbidden high-level mutations}. Direct ablations reproduce the expected mechanism separations, and an independent implementation recomputes every score and analysis field without disagreement.

The post-saturation P1-X study extends the same principle from hidden formulation shifts to heterogeneous scientific revision responsibility. In the prospectively repaired, comparator-fair V2, P1-X achieves \textbf{400/400 exact revision decisions} versus 275/400 for a donor-complete greedy product, a paired advantage of 0.3125 with 95\% CI [0.2675,0.3575]. It preserves zero false high-level reframes and zero protected-invariant violations across the exact domains, and an independent implementation reproduces the protected identities, result hashes, and headline counts.

An ideal donor product supplied the same responsibility, minimality, preservation, reopening, and authority semantics also reaches 400/400. This is a portability result: the scientific escalation contract is sufficient whether implemented centrally or as a correctly integrated donor product. ORION's advantage over the realistic donor-complete comparator comes from making that coupling explicit, complete, and enforceable.

The widest supported conclusion is therefore not merely that one system wins a hidden-shift benchmark. It is a scientific escalation law over the registered families: \textbf{failure is evidence for diagnosis, not authority for reformulation; high-level scientific change becomes admissible only when protected necessity evidence rules out lower-level repair}. Under that law, ORION is both dramatically more successful when the formulation is wrong and perfectly selective when lower-level repair is sufficient. The negative historical study supplies the falsifier; the powered primary/replication and P1-X supply the prospective repaired evidence. Model-general and open-ended deployment are the natural next transfer tests of an already demonstrated contract rather than caveats on the protected result.
